% ============================================================
%  Chapter 5 — Discussion
% ============================================================
\chapter{Discussion}
\label{ch:discussion}

The results of Chapter~\ref{ch:results} raise two distinct questions that are
easily conflated. The first is technical: did the machine-learning models do
what they were asked to do? The second is substantive: is the resulting
groundwater potential map actually useful and correct? Keeping these apart is
essential to an honest reading of this work, and the discussion below is
organized around that distinction---beginning with model behaviour and feature
interpretation, then confronting the label-circularity problem directly, and
finally drawing out what the validation can and cannot establish, how the map
compares with regional studies, where it is likely to be wrong, and what it is
good for.

\section{Model Performance in Context}
\label{sec:discuss_performance}

XGBoost achieved the best results across every metric, though only narrowly ahead
of LightGBM and the voting ensemble. The closeness of the three boosting-based
models is itself informative: when several well-regularized learners converge on
almost identical performance, it usually means they have extracted essentially
all of the available signal, and the remaining differences reflect minor
implementation details rather than any deep advantage. XGBoost's slight edge is
consistent with its regularized objective and column subsampling
\cite{chen2016xgboost}, which tend to help on moderate-dimensional inputs like the
nine-feature stack used here. The ordering---XGBoost first, random forest a clear
step behind---also matches what Sahin \cite{sahin2020ensembletree} found when
comparing the same families for landslide susceptibility, where boosting
outperformed bagging on structurally similar spatial data. That an independent
study on a different hazard reproduces the same ranking gives some confidence that
the result reflects a real property of the methods rather than a quirk of this
dataset.

The absolute numbers, however, demand caution rather than celebration. An accuracy
near 99\% and an AUC of 0.9998 would be extraordinary for a model validated
against field observations; here they are unsurprising, because the target is a
deterministic function of the input features. The next section takes this point
head-on, since it governs how every performance figure in this thesis should be
read.

\section{Feature Importance and Hydrogeology}
\label{sec:discuss_features}

Both the correlation analysis and the trained models identify soil texture and
rainfall as the two controlling factors, and the agreement between a simple linear
diagnostic and a nonlinear ensemble is reassuring. Soil texture's dominance
(XGBoost importance 0.831; $r = 0.836$) is hydrogeologically sensible: texture is
the surface expression of permeability, and permeability decides whether rainfall
infiltrates to recharge an aquifer or runs off. This is also consistent with the
reference study \cite{waqas2022chakwal}, in which geology and land use---the
nearest proxies for subsurface permeability---together carried the largest weight,
and it lends after-the-fact support to the decision to source soil texture from
SoilGrids \cite{poggio2021soilgrids} rather than omit it as the reference study
effectively did.

Rainfall's second place is equally expected in a monsoon-fed, recharge-limited
setting, and it echoes Khan et al. \cite{khan2023rainfallrunoff}, who found
rainfall to be the dominant control on the hydrological response of these same
basins. The near-zero importance of slope ($r = 0.051$) deserves a brief comment
because it might seem to contradict the GWPM literature, where slope is usually a
meaningful factor. The explanation is regional: across a comparatively subdued
plateau, slope varies too little to separate groundwater zones, whereas in the
mountainous catchments where slope is informative the relief is far greater. The
finding is therefore specific to this terrain and should not be generalized.

\section{The Label-Circularity Problem}
\label{sec:discuss_circularity}

The central methodological issue in this work is that the training labels were
generated from the same nine features that the models take as input. The
weighted overlay computes a score from the features, thresholds it into three
classes, and those classes become the targets the models learn to predict. The
consequence is a form of circularity: the models are, in effect, learning to
reproduce a fixed arithmetic combination of their own inputs. The 99.05\%
agreement between the XGBoost map and the overlay labels
(Section~\ref{sec:agreement}), and the near-perfect test metrics, are therefore
largely a measure of how completely the model recovered that combination---not
independent evidence that the map matches real groundwater.

It would be dishonest to present the high accuracy as though it validated the map
against reality, and this thesis does not. What the accuracy does establish is
worth stating precisely, because it is not nothing. It shows that a single learned
model can reproduce the weighted overlay across 37~million pixels faithfully and
efficiently, which means the overlay can be replaced by a model that, unlike the
overlay, does not require fixed expert weights, can represent nonlinear
interactions among features, and yields a continuous probability surface rather
than a hard class. In other words, the contribution of the ML step is not a more
accurate map than the overlay---by construction it is almost the same map---but a
more flexible, reproducible, and extensible way of producing it. This is the same
logic under which knowledge-driven labels have been used to train spatial models
elsewhere; Roy and Saha \cite{roy2019knowledgedriven}, for instance, trained on
expert-weighted targets and treated the exercise as legitimate precisely because
the weighted product encodes real domain knowledge.

The honest conclusion is that internal accuracy cannot certify the map, and a
different kind of evidence is needed. That is the purpose of the spatial-realism
validation, discussed next.

\section{What the Validation Establishes---and What It Does Not}
\label{sec:discuss_validation}

The four checks in Section~\ref{sec:validation_results} were designed to test the
map against evidence that played no part in its construction, and on their own
terms they are encouraging. The map places all four known alluvial sites in
Medium or High ground; it separates independent rain-gauge stations by
hydrogeological setting at a statistically significant level (Mann--Whitney
$p = 0.023$); its 20\% High fraction matches the magnitude of an independent
regional AHP study; and it reproduces the expected district gradient perfectly
(Spearman $\rho = 1.00$). Crucially, none of these uses the overlay labels, so
together they speak to the substantive question that internal accuracy cannot.

Their limits should be stated just as plainly. The known-sites check rests on only
four points and is qualitative. The rain-gauge test, though significant, is based
on a modest number of stations. The cross-study comparison establishes consistency
of magnitude, not agreement pixel by pixel. And the district gradient, while the
strongest of the four because it aggregates large populations, confirms only that
the broad north--south ordering is right, not that any particular hectare is
correctly classified. What the validation supports, then, is a measured claim: the
map is spatially realistic---it is right in its large-scale structure and on the
independent points available---rather than a claim that it is accurate everywhere.
For a data-scarce region this is, realistically, the strongest validation the
available evidence will bear, and it is a more honest position than leaning on the
99\% figure. The recurring difficulty of obtaining adequate ground truth for
groundwater models is itself well documented \cite{hai2022gwlreview}; this thesis
does not escape it, but it does try to be candid about it.

\section{Comparison with the Reference and Regional Studies}
\label{sec:discuss_comparison}

Set against the reference study of Waqas et al. \cite{waqas2022chakwal}, the
advance here is one of scope and method rather than of validated accuracy
(Table~\ref{tab:comparison}). The present work covers all five districts of the
plateau instead of one, uses nine features instead of five, replaces manual
overlay with a learned model, and reports quantitative performance and an
independent validation where the reference study reported neither. These are real
improvements in reproducibility and coverage, and they are the substantive
contribution; the headline accuracy, as discussed, is not directly comparable
because the reference study did not train on its own overlay.

\begin{table}[H]
\centering
\caption{Comparison of this thesis with the reference study of Waqas et al.
         \cite{waqas2022chakwal}.}
\label{tab:comparison}
\small
\begin{tabular}{lll}
\toprule
Aspect & Waqas et al.\ (2022) & This thesis \\
\midrule
Study area & Chakwal (1 district) & Full plateau (5 districts) \\
Method     & Manual weighted overlay & ML (RF, XGBoost, LightGBM, ensemble) \\
Features   & 5 & 9 \\
Validation & Secondary data only & 5-fold CV + spatial-realism checks \\
Output     & Static maps & GeoTIFF class + probability rasters \\
Resolution & Not specified & 30~m \\
\bottomrule
\end{tabular}
\end{table}

The map also sits comfortably alongside other recent work on the region. Hafeez et
al. \cite{hafeez2025rwh}, mapping rainwater-harvesting suitability across
essentially the same districts by AHP, identify a similar minority of the area as
most favourable and locate it in the same alluvial tracts, which is the
cross-study consistency reported in Section~\ref{sec:validation_results}. The
broader move from knowledge-driven overlay to data-driven prediction that this
thesis carries out for groundwater has already been found beneficial for related
problems---Arabameri et al. \cite{arabameri2018shahroud}, comparing approaches on
a single plain, reported data-driven methods matching or exceeding multi-criteria
ones---so the methodological direction taken here is well supported even if the
circularity caveat tempers the accuracy gain.

\section{Likely Errors: Hard-Rock Terrain and the Missing Lithology Layer}
\label{sec:discuss_hardrock}

The clearest weakness of the map is expected to lie in the hard-rock uplands. The
nine features capture surface and climatic conditions well, but they say nothing
direct about subsurface lithology or fracturing, and in crystalline or massively
bedded terrain these are precisely the controls that matter. Benjmel et al.
\cite{benjmel2020morocco} showed that in such settings lineament density and
geological structure dominate groundwater occurrence, with fracture zones acting
as conduits through otherwise impermeable rock. Where parts of the Margalla Hills
or the Kala Chitta Range carry high rainfall and dense vegetation, the model is
liable to read those surface signals as high potential even though the underlying
rock may store little water. The absence of a lithological or lineament layer is
thus the most probable source of systematic over-prediction, and adding a
digitized geological map---from Geological Survey of Pakistan sheets---as a tenth
feature is the single most valuable improvement available to future work.

A second, subtler caveat concerns the distinction between potential and
potability. The model predicts where groundwater is likely to occur, not whether
it is fit to use. In the peri-urban Rawalpindi--Islamabad belt, shallow aquifers
carry a documented contamination risk \cite{waseem2024waterquality}, so a
high-potential classification there should be read as "water is likely present,"
not "water is safe to drink." A water-quality layer would be needed to bridge that
gap.

\section{Practical Implications}
\label{sec:discuss_implications}

With those caveats in mind, the map still carries practical value for an area
where groundwater siting is currently done with little spatial guidance. The High
potential zone---about 20\% of the plateau, on the order of a few thousand square
kilometres---is concentrated along the Soan basin and the northern and riverine
alluvial corridors, and these are the natural first places to consider for new
wells or managed recharge. The Low zone, dominating the southern hard-rock
terrain, flags where deep borehole investment is least likely to pay off, subject
to the lithology caveat above. Because the outputs are standard GeoTIFF rasters at
30~m resolution---both a hard classification and a continuous probability
surface---they can be brought directly into the GIS workflows of agencies such as
PCRWR and the provincial irrigation department for district-level planning. In a
region where rain-fed agriculture leaves households acutely dependent on
groundwater through the dry season \cite{sarwar2022soandrought}, even a
spatially-realistic first-order map is more useful than none.

\section{Limitations}
\label{sec:discuss_limitations}

Several limitations have been noted in passing and are collected here for clarity.
The foremost is label circularity: with no plateau-wide field measurements, the
models were trained on knowledge-based labels, so internal accuracy measures
fidelity to the overlay rather than to ground truth. The validation, while
independent, rests on relatively few points and establishes spatial realism rather
than pixel-level accuracy. The feature set omits subsurface lithology and
lineaments, which likely causes over-prediction in hard-rock zones, and it
includes no water-quality information. Finally, the static input layers---a
multi-year rainfall sum, single-period spectral indices---describe an average
condition and do not capture seasonal or interannual variation in groundwater.
None of these undermines the central contribution, but each marks a place where
the map should be used with appropriate caution and where future work can
strengthen it.
